\section{Riscos e Mitigacoes}

\subsection{Analise de Riscos}

\subsubsection{Matriz de Probabilidade e Impacto}

\begin{longtable}{|p{4cm}|p{2cm}|p{2cm}|p{5cm}|}
\hline
\textbf{Risco} & \textbf{Prob.} & \textbf{Impacto} & \textbf{Descricao} \\
\hline
\endhead

Seguranca de Dados & Media & Alto & Vazamento de informacoes financeiras sensiveis \\
\hline

Atraso no Desenvolvimento & Alta & Medio & Cronograma nao cumprido por complexidade tecnica \\
\hline

Problemas de Performance & Media & Medio & Sistema lento com muitos usuarios \\
\hline

Rejeicao nas App Stores & Baixa & Alto & Apps rejeitados por politicas das stores \\
\hline

Mudancas Regulatorias & Baixa & Alto & Novas leis sobre dados financeiros \\
\hline

Concorrencia & Alta & Medio & Lancamento de produtos similares \\
\hline

Problemas de UX & Media & Alto & Interface confusa ou dificil de usar \\
\hline

Falhas de Infraestrutura & Baixa & Alto & Indisponibilidade do sistema \\
\hline

\end{longtable}

\subsection{Riscos Tecnicos}

\subsubsection{RT001 - Seguranca de Dados Financeiros}

\paragraph{Descricao}
Dados financeiros sao extremamente sensiveis e qualquer vazamento pode causar graves problemas legais e de reputacao.

\paragraph{Probabilidade}: Media (30\%)
\paragraph{Impacto}: Alto

\paragraph{Fatores de Risco}
\begin{itemize}
    \item Ataques ciberneticos direcionados
    \item Vulnerabilidades no codigo
    \item Configuracoes incorretas de seguranca
    \item Acesso nao autorizado por funcionarios
\end{itemize}

\paragraph{Estrategias de Mitigacao}
\begin{itemize}
    \item \textbf{Preventivas}:
    \begin{itemize}
        \item Criptografia AES-256 para dados sensiveis
        \item HTTPS obrigatorio em todas as comunicacoes
        \item Autenticacao multi-fator
        \item Auditoria de codigo por especialistas em seguranca
        \item Penetration testing regular
    \end{itemize}
    \item \textbf{Detectivas}:
    \begin{itemize}
        \item Monitoramento de acessos suspeitos
        \item Logs detalhados de todas as operacoes
        \item Alertas automaticos de tentativas de invasao
    \end{itemize}
    \item \textbf{Corretivas}:
    \begin{itemize}
        \item Plano de resposta a incidentes
        \item Backup criptografado e testado
        \item Protocolo de comunicacao com usuarios
    \end{itemize}
\end{itemize}

\subsubsection{RT002 - Performance e Escalabilidade}

\paragraph{Descricao}
Sistema pode nao suportar o crescimento de usuarios ou apresentar lentidao.

\paragraph{Probabilidade}: Media (40\%)
\paragraph{Impacto}: Medio

\paragraph{Estrategias de Mitigacao}
\begin{itemize}
    \item Arquitetura escalavel desde o inicio
    \item Cache Redis para operacoes frequentes
    \item CDN para assets estaticos
    \item Monitoramento de performance em tempo real
    \item Testes de carga automatizados
    \item Auto-scaling na infraestrutura cloud
\end{itemize}

\subsubsection{RT003 - Complexidade do Desenvolvimento}

\paragraph{Descricao}
Subestimacao da complexidade pode levar a atrasos significativos.

\paragraph{Probabilidade}: Alta (60\%)
\paragraph{Impacto}: Medio

\paragraph{Estrategias de Mitigacao}
\begin{itemize}
    \item Desenvolvimento iterativo com entregas frequentes
    \item Prototipagem antes da implementacao
    \item Code review obrigatorio
    \item Documentacao tecnica detalhada
    \item Buffer de 20\% no cronograma
    \item Equipe com experiencia nas tecnologias escolhidas
\end{itemize}

\subsection{Riscos de Negocio}

\subsubsection{RN001 - Aceitacao do Usuario}

\paragraph{Descricao}
Usuarios podem nao adotar a metodologia proposta ou achar a interface complexa.

\paragraph{Probabilidade}: Media (35\%)
\paragraph{Impacto}: Alto

\paragraph{Estrategias de Mitigacao}
\begin{itemize}
    \item Pesquisa de mercado detalhada
    \item Testes de usabilidade com usuarios reais
    \item MVP com funcionalidades essenciais
    \item Onboarding guiado e intuitivo
    \item Feedback continuo dos usuarios
    \item Tutorial interativo e help context-sensitive
\end{itemize}

\subsubsection{RN002 - Concorrencia}

\paragraph{Descricao}
Entrada de competidores com recursos superiores ou lancamento antes do previsto.

\paragraph{Probabilidade}: Alta (70\%)
\paragraph{Impacto}: Medio

\paragraph{Estrategias de Mitigacao}
\begin{itemize}
    \item Foco na diferenciacao da metodologia unica
    \item Lancamento rapido do MVP
    \item Construcao de comunidade de usuarios
    \item Melhoria continua baseada em feedback
    \item Parcerias estrategicas
    \item Marketing focado nos beneficios unicos
\end{itemize}

\subsection{Riscos Regulatorios}

\subsubsection{RR001 - LGPD e Protecao de Dados}

\paragraph{Descricao}
Mudancas na legislacao de protecao de dados podem exigir alteracoes no sistema.

\paragraph{Probabilidade}: Baixa (20\%)
\paragraph{Impacto}: Alto

\paragraph{Estrategias de Mitigacao}
\begin{itemize}
    \item Conformidade total com LGPD desde o inicio
    \item Consulta juridica especializada
    \item Termo de consentimento claro e especifico
    \item Funcionalidades de exclusao de dados
    \item Portabilidade de dados do usuario
    \item Acompanhamento continuo de mudancas legais
\end{itemize}

\subsubsection{RR002 - Regulamentacao de Fintechs}

\paragraph{Descricao}
Novas regulamentacoes podem afetar aplicacoes financeiras.

\paragraph{Probabilidade}: Baixa (15\%)
\paragraph{Impacto}: Medio

\paragraph{Estrategias de Mitigacao}
\begin{itemize}
    \item Foco em ferramenta de organizacao, nao transacoes
    \item Nao custodia de recursos financeiros
    \item Relacionamento com associacoes do setor
    \item Monitoramento de mudancas regulatorias
\end{itemize}

\subsection{Riscos Operacionais}

\subsubsection{RO001 - Disponibilidade da Equipe}

\paragraph{Descricao}
Membros da equipe podem nao estar disponiveis durante periodos criticos.

\paragraph{Probabilidade}: Media (30\%)
\paragraph{Impacto}: Medio

\paragraph{Estrategias de Mitigacao}
\begin{itemize}
    \item Documentacao tecnica completa
    \item Knowledge sharing regular entre a equipe
    \item Redundancia nas competencias criticas
    \item Cronograma flexivel com buffers
    \item Backup de fornecedores/freelancers
\end{itemize}

\subsubsection{RO002 - Falhas de Infraestrutura}

\paragraph{Descricao}
Problemas com provedores de cloud ou servicos terceiros.

\paragraph{Probabilidade}: Baixa (15\%)
\paragraph{Impacto}: Alto

\paragraph{Estrategias de Mitigacao}
\begin{itemize}
    \item Multiplas zonas de disponibilidade
    \item Backup automatico e testado regularmente
    \item SLA definido com provedores
    \item Monitoramento 24/7
    \item Plano de disaster recovery
    \item Redundancia em servicos criticos
\end{itemize}

\subsection{Plano de Contingencia}

\subsubsection{Cenario 1: Atraso Significativo (>4 semanas)}
\begin{itemize}
    \item Reducao de escopo priorizando features essenciais
    \item Aumento temporario da equipe
    \item Extensao do cronograma com comunicacao transparente
    \item Lancamento em fases (soft launch)
\end{itemize}

\subsubsection{Cenario 2: Problemas Graves de Seguranca}
\begin{itemize}
    \item Suspensao imediata do servico se necessario
    \item Comunicacao transparente com usuarios
    \item Auditoria externa de seguranca
    \item Implementacao de correcoes antes do retorno
\end{itemize}

\subsubsection{Cenario 3: Rejeicao Nas App Stores}
\begin{itemize}
    \item Analise detalhada dos motivos da rejeicao
    \item Correcoes prioritarias
    \item Foco na versao web enquanto resolve o mobile
    \item Submissao em lotes menores se necessario
\end{itemize}

\subsection{Monitoramento de Riscos}

\subsubsection{Indicadores de Alerta}
\begin{itemize}
    \item Atraso >1 semana em qualquer sprint
    \item Taxa de bugs criticos >5\% dos testes
    \item Performance >3 segundos para operacoes basicas
    \item Feedback negativo >30\% nos testes de usuario
    \item Violacao de SLA de infraestrutura
\end{itemize}

\subsubsection{Revisoes Regulares}
\begin{itemize}
    \item Avaliacao semanal de riscos em sprint reviews
    \item Revisao mensal de riscos estrategicos
    \item Atualizacao trimestral da matriz de riscos
    \item Post-mortem apos qualquer incidente
\end{itemize}
